\subsection{Instruction Payload Ordering}
The instruction header occupies one byte for extrashort forms, two bytes for short forms, and at least two start
bytes for extended forms. If a descriptor, immediate, displacement, or bitmap follows, those payload bytes
occupy increasing addresses in decode order.

Signed displacement and immediate payloads use two's-complement representation at their declared width before any
sign extension performed by the consuming instruction or EA form.

\manualtablecaption{Immediate, Displacement, and Address Payload Order}
\begin{manuallongtable}{@{}>{\raggedright\arraybackslash}p{0.75in}>{\raggedright\arraybackslash}p{0.45in}>{\raggedright\arraybackslash}p{4.20in}@{}}
\toprule
\rowcolor{ManualHeaderFill}
\textbf{Payload} & \textbf{Bytes} & \textbf{Memory Order}\\
\midrule
\endfirsthead
\multicolumn{3}{l}{\scriptsize\itshape Table \themanualtable\ (continued)}\\
\toprule
\rowcolor{ManualHeaderFill}
\textbf{Payload} & \textbf{Bytes} & \textbf{Memory Order}\\
\midrule
\endhead
\texttt{imm8} & 1 & single byte\\
\texttt{imm16} & 2 & bits 7..0 at the lower byte, then bits 15..8\\
\texttt{imm32} & 4 & least significant byte first\\
\texttt{imm64} & 8 & least significant byte first\\
\texttt{disp8} & 1 & single signed two's-complement byte\\
\texttt{disp16} & 2 & signed two's-complement value, little-endian\\
\texttt{disp32} & 4 & signed value, least significant byte first\\
\texttt{disp64} & 8 & declared-width payload, least significant byte first\\
\texttt{abs32} & 4 & absolute-address payload, least significant byte first\\
\texttt{abs64} & 8 & absolute-address payload, least significant byte first\\
\bottomrule
\end{manuallongtable}

\manualtablecaption{Immediate Operand Interpretation}
\begin{manuallongtable}{@{}>{\raggedright\arraybackslash}p{0.95in}>{\raggedright\arraybackslash}p{0.38in}>{\raggedright\arraybackslash}p{0.62in}>{\raggedright\arraybackslash}p{0.62in}>{\raggedright\arraybackslash}p{0.90in}>{\raggedright\arraybackslash}p{1.70in}@{}}
\toprule
\rowcolor{ManualHeaderFill}
\textbf{Operand} & \textbf{Bits} & \textbf{Value} & \textbf{Range} & \textbf{Extension} & \textbf{Applies When}\\
\midrule
\endfirsthead
\multicolumn{6}{l}{\scriptsize\itshape Table \themanualtable\ (continued)}\\
\toprule
\rowcolor{ManualHeaderFill}
\textbf{Operand} & \textbf{Bits} & \textbf{Value} & \textbf{Range} & \textbf{Extension} & \textbf{Applies When}\\
\midrule
\endhead
\texttt{pair\_id} & @PAIR_ID_WIDTH@ & identifier & @PAIR_ID_RANGE@ & none & PUSHP/POPP canonical register-pair selector\\
\texttt{pt\_level} & @PT_LEVEL_WIDTH@ & enumeration & @PT_LEVEL_RANGE@ & none & PTQUERY page-table level selector\\
\texttt{flags\_bitmap} & @FLAGS_BITMAP_WIDTH@ & bitmap & @FLAGS_BITMAP_RANGE@ & none & SETF/CLRF integer FLAGS masks\\
\texttt{imm6} & @IMM6_WIDTH@ & unsigned & @IMM6_RANGE@ & zero-extend & integer forms with an explicit B/W/L/Q operation size\\
\texttt{imm7} & @IMM7_WIDTH@ & unsigned & @IMM7_RANGE@ & zero-extend & EXTRACT register-pair forms\\
\bottomrule
\end{manuallongtable}

PUSHP and POPP interpret \texttt{pair\_id} as a canonical register-pair index, so every encoding from 0 through 7 is valid.
PTQUERY accepts \texttt{pt\_level} values 1 through 5; encodings 0, 6, and 7 are reserved.
SETF and CLRF interpret \texttt{flags\_bitmap} as the low FLAGS nibble: bit 3 selects Z, bit 2 selects N, bit 1 selects C,
and bit 0 selects V.

An \texttt{imm6} value is zero-extended to the selected operation size before use when that size is wider than six
bits. EXTRACT uses \texttt{imm7} as the offset into its concatenated register pair.
